\chapter*{Заключение}
\addcontentsline{toc}{chapter}{Заключение}

В работе были исследованы методы прогнозирования осадков и детекции аномальных осадочных событий на климатических данных. Основная идея состояла в том, чтобы рассматривать осадки не только как непрерывную величину, но и как источник редких событий, определяемых через превышение локального высококвантильного порога.

В станционной части была построена полная процедура обработки данных RusWeather-GF: формирование лаговых и агрегированных признаков, временное разбиение выборки с учетом защитного интервала, обучение климатологии, линейной модели и градиентных бустингов. Результаты показали, что климатология является сильной базовой моделью, линейная модель недостаточно хорошо описывает нелинейную структуру осадков, а лучшие результаты среди табличных методов дают LightGBM и XGBoost. На горизонтах 24--120 часов XGBoost показывает наилучшие событийные метрики, тогда как на горизонте 240 часов климатология становится особенно сильным ориентиром.

Во второй части была рассмотрена региональная нейросетевая модель Earthformer-M для сеточных полей. Модель использует многоканальные последовательности метеорологических полей, отдельную обработку динамических и статических признаков, 3D-сверточный входной тракт, cuboid transformer-блоки, temporal attention pooling и условную выходную часть, зависящую от горизонта прогноза и номера ансамблевого члена. Обучение проводится в многозадачной постановке: одновременно оптимизируются регрессионные компоненты по осадкам, квантильная компонента ансамбля, soft-event loss, Brier-компонента и вспомогательная событийная голова.

По результатам обучения и проверки на валидационном периоде можно выделить несколько содержательных наблюдений. Во-первых, функция потерь устойчиво уменьшалась, а якорная метрика ROCSS выходила на плато после первых эпох, что показывает необходимость выбирать модель не только по общей функции потерь, но и по метрикам редких событий. Во-вторых, региональная модель показывает очень сильное качество ранжирования на горизонте 24 часа, но на средних и дальних горизонтах качество заметно снижается. В-третьих, новый запуск выявил существенное завышение прогнозных вероятностей и положительное смещение по сумме осадков, поэтому вероятностный выход требует отдельной калибровки перед практическим использованием.

Ключевой результат работы состоит в том, что детекция аномалий в климатических данных может быть встроена в задачу прогноза через квантильное определение события. В станционной постановке это делается через сравнение непрерывного прогноза с локальным порогом станции, а в сеточной постановке — через локальную карту порогов и вероятностную оценку превышения. Такая схема сохраняет физическую интерпретацию целевой величины и одновременно дает инструмент для оценки риска экстремальных осадков.

В ходе работы были решены поставленные задачи: сформулирована событийная постановка, подготовлены станционные признаки, обучены и сравнены базовые модели, описана региональная нейросетевая модель Earthformer-M, рассчитаны вероятностные метрики, построены графики калибровки, пространственных примеров, временных смещений и ансамблевой неопределенности.

Практическое значение работы заключается в том, что предложенная структура эксперимента может использоваться при разработке систем прогноза осадков на станциях и при построении пространственно согласованных систем оценки риска экстремальных осадков. Дальнейшее развитие работы может быть связано с более строгой калибровкой вероятностей, расширением набора регионов, доработкой функции потерь для уменьшения положительного смещения осадков и переносом подхода на другие типы климатических аномалий: температурные экстремумы, засухи и комбинированные гидроклиматические события.
